% sign_schur.tex -- round 394, lemma-first attack on the
% sign binding of the Uvarov image K^mu[Xi] (the Assist rest
% named by r392).
% Outcome: REDUCED.  Builds with pdflatex.
% Companion machine check: verify_sign_schur.py.
% Discovery census: experiments/tfpt-discovery/sign_schur_probe.py.
% NO RH CLAIM.  Research documentation, not a theorem of RH.
\documentclass[11pt]{article}

\usepackage[a4paper,margin=2.7cm]{geometry}
\usepackage{amsmath,amssymb,amsthm}
\usepackage{booktabs}
\usepackage{microtype}
\usepackage{xcolor}
\usepackage[colorlinks=true,linkcolor=blue!50!black,urlcolor=blue!50!black]{hyperref}

\newtheorem{lemma}{Lemma}
\newtheorem{prop}{Proposition}
\newtheorem{definition}{Definition}
\theoremstyle{remark}
\newtheorem{remark}{Remark}

\newcommand{\R}{\mathbb{R}}
\newcommand{\N}{\mathbb{N}}

\title{\bfseries Sign-Schur, reduced\\[2mm]
\large Dirichlet zones of the Christoffel--Darboux kernel,
not a bipartition of the $\nu$-nodes}
\author{Stefan Hamann}
\date{August 28, 2026}

\begin{document}
\maketitle

\begin{abstract}\noindent
Round 392 located the $99.7\%$ Assist cancellation in the
\emph{signs} of the Uvarov image $K^\mu[\Xi]$, not in the
positive $\tau$-quotients of occupied Fej\'er.
The candidate mechanism was a source-defined checkerboard
along the ordered $\nu$-nodes, inherited from the oscillating
Chebyshev--Christoffel--Darboux kernel, which via a diagonal
sign conjugation $D_\sigma E D_\sigma$ would force
$\lambda_{\max}(E_k)\le\mathrm{maxdiag}\cdot(1+o(1))$.
What is proved: on the cosine grid the Chebyshev-$T$ kernel
is the Dirichlet combination of~\cite{r387}, the sign of the
difference term
$D_{n-1}(\theta_i-\theta_j)$ is zonal
(depends only on which $\pi$-interval $(n-\tfrac12)|\Delta\theta|$
occupies), and
$\lvert D_n(\alpha)\rvert\le\min(2n+1,1/\lvert\sin(\alpha/2)\rvert)$.
The hoped M-matrix bound is false even in $2\times 2$:
$\bigl(\begin{smallmatrix}1&-c\\-c&1\end{smallmatrix}\bigr)$
has $\lambda=1+c>\mathrm{maxdiag}$.
What is \emph{not} proved: rank-1 conjugability.
Mass-weighted checkerboard agreement is a coin ($0.504$ on
FRAME-A $w=9$; core-$42$ in $[0.429,0.512]$).
The large couplings \emph{do} inherit Chebyshev signs
(mass-weighted $0.810$ on FRAME-A, $[0.810,0.894]$ on core-$42$),
but the $11$--$19\%$ residual plus the failed bipartition block
every source-defined $k^*>2N/5$.
Two-period lag-$1$ is in-phase (named kill).
Scramble holds the mass-map and dies on the envelope.
No RH claim.
\end{abstract}

\medskip\noindent
\textbf{Status.}\enspace
Lemmas~\ref{lem:dirichlet-sign}--\ref{lem:zfalse} are proved
(finite identities / a named counterexample).
Lemma~\ref{lem:inphase} is a machine identity on the two-period
mesh.
Propositions~\ref{prop:map}--\ref{prop:red} are a finite census
and a named reduction.
No RH claim.  No $L^*$ claim.

\section{Recollection}\label{sec:rec}

Write $E_{ij}=\sqrt{v_i v_j}\,K_n^\mu(y_i,y_j)$ for the
IIKS Gram at the $\nu$-nodes (the Uvarov image of~\cite{r392}).
On FRAME-A $w=9$, $k=183$:
$\lambda_{\max}=0.999830$, $\mathrm{maxdiag}=0.9614$,
$\mathrm{assist}=0.0399$, relative Gershgorin $g_A=13.32$,
cancellation $0.9970$.
Gershgorin of the absolute matrix is too crude~\cite{r387}.
The remaining question is whether the off-diagonal
\emph{sign pattern} is structural enough to force
$\lambda_{\max}\le\mathrm{maxdiag}\cdot(1+o(1))$ from the
source (F1-runs, separation, $\rho_{\mathrm{AP}}<1/5$).

The candidate: order the $\nu$-nodes by angle and hope
$\operatorname{sign}(K_{ij})=(-1)^{i-j}$ (checkerboard),
equivalently a sign vector $\sigma$ with
$\sigma_i\sigma_j K_{ij}\le 0$ for almost all pairs.
Then $D_\sigma E D_\sigma$ would be a Z-matrix, and
M-matrix / Perron theory would bound $\lambda_{\max}$.

\section{Dirichlet signs on the cosine grid}\label{sec:dirichlet}

The Chebyshev-$T$ Christoffel--Darboux kernel is~\cite{r387}
\[
K_n^T(\theta,\phi)
=
\frac{D_{n-1}(\theta-\phi)+D_{n-1}(\theta+\phi)}{2\pi},
\qquad
D_m(\alpha)
=
\frac{\sin\bigl((m+\tfrac12)\alpha\bigr)}{\sin(\alpha/2)}.
\]

\begin{lemma}[Dirichlet sign]\label{lem:dirichlet-sign}
For $\alpha\not\in 2\pi\mathbb{Z}$,
\[
\operatorname{sign}(D_m(\alpha))
=
\operatorname{sign}\bigl(\sin((m+\tfrac12)\alpha)\bigr)
\cdot
\operatorname{sign}\bigl(\sin(\alpha/2)\bigr).
\]
On the uniform cosine grid $\theta_j=\pi j/S$ the difference
term is zonal: its sign depends only on which $\pi$-interval
contains $(m+\tfrac12)\lvert\theta_i-\theta_j\rvert$.
Machine identity on $S=16,12,21$.
\end{lemma}

\begin{lemma}[Dirichlet envelope]\label{lem:envelope}
$\lvert D_n(\alpha)\rvert\le\min\bigl(2n+1,\,1/\lvert\sin(\alpha/2)\rvert\bigr)$
for $\alpha\not\in 2\pi\mathbb{Z}$.
Equality of the pointwise ratio to $1$ is attained on the
tested mesh; the bound is classical.
\end{lemma}

The reflected sum $D_{n-1}(\theta_i+\theta_j)$ flips about
one eighth of the pair-signs on a small grid (agree $0.875$
against the difference term alone).
The zonal law is a theorem of the Chebyshev kernel, not of
the $\mu$-OP Gram.

\section{The hoped bound is false in \texorpdfstring{$2\times 2$}{2x2}}\label{sec:z}

\begin{lemma}[Z-matrix bound false]\label{lem:zfalse}
For $0<c<1$ the matrix
$\bigl(\begin{smallmatrix}1&-c\\-c&1\end{smallmatrix}\bigr)$
is positive definite with negative off-diagonals and
$\lambda_{\max}=1+c>\mathrm{maxdiag}$.
A diagonal sign conjugation to a Z-matrix does \emph{not}
imply $\lambda_{\max}\le\mathrm{maxdiag}$.
\end{lemma}

The $2\times 2$ already kills the M-matrix/Perron slogan as a
closing deduction, even granting perfect conjugacy.
What remains is a \emph{quantitative} remainder: if the
positive off-diagonal mass after a \emph{source-defined}
conjugation is $o(\mathrm{maxdiag})$, Weyl's inequality still
needs a bound on $\lambda_{\max}$ of that remainder.
Oracle conjugation by the true top eigenvector of $E$ on
FRAME-A leaves a positive-remainder row-sum $1.90>1$.

\section{The sign map, measured}\label{sec:map}

\begin{prop}[sign census]\label{prop:map}
On FRAME-A $w=9$ at $k=183$, ordering $\nu$-nodes by angle:
\begin{itemize}
\item checkerboard mass-weighted agreement $0.504$ (coin);
\item Chebyshev-CD mass-weighted agreement $0.810$,
pair-count $0.721$ (small entries flip);
\item nearest-neighbour negative mass-share $0.468$
(not an M-matrix path);
\item last twelve Chebyshev mass-agreements in
$[0.801,0.829]$ (robust in $k$, not a SATZ);
\item at the r382 entry $k=\lfloor 2N/5\rfloor=73$ the
nearest neighbours are \emph{constructive}
(negative mass-share $0.039$).
\end{itemize}
Core-$42$ at the wall: Chebyshev mass-agreement in
$[0.810,0.894]$, checkerboard in $[0.429,0.512]$,
nearest-neighbour negative share in $[0.349,0.543]$,
$42/42$ with $\lambda_{\max}<1$.
Builder fallback: core-$42$ only.
\end{prop}

The $\tau$-deformation of r392 is small in the
\emph{formula} for $\gamma$, but not small enough to transfer
the Dirichlet-zone SATZ to $K^\mu$: eleven to nineteen percent
of the off-diagonal \emph{mass} disagrees, and the pair-count
disagrees at one quarter to one third.
The large couplings follow Chebyshev; the bipartition does not.

\section{Adversaries}\label{sec:kills}

\begin{lemma}[two-period in-phase]\label{lem:inphase}
On the two-period mesh $S=81$, $c=\tfrac23$, at depth $k=22$
every nearest-neighbour off-diagonal of $E$ is positive
($40/0$) and $\lambda_{\max}=1.0288>1$.
A checkerboard conjugation along the chain is impossible:
the short-lag couplings are in-phase.
At $k=40$ the Chebyshev (and $\mu$-CD) sign pattern
\emph{is} exactly the anti-checkerboard, but too late:
$\mathrm{maxdiag}=1.317$ already exceeds $1$.
\end{lemma}

\paragraph{Named kills.}
Scramble occupation seed $1$ (the r388 $\mathrm{d}\Delta=2.16\cdot10^6$
adversary) keeps Chebyshev mass-agreement $0.788$ and diverges
at $\lambda_{\max}=9.17\cdot10^6$: it breaks the
\emph{envelope} (occupation / maxdiag), not the mass-weighted
sign map.
The mutant $E\mapsto\lvert E\rvert$ (signs dropped) yields
$\lambda_{\max}(\lvert E\rvert)=1.683>1$ on FRAME-A; the
relative $g_A=13.32$ is sign-blind and unchanged.
Assist $0.0399$ is entirely the signed remainder.

\section{What remains}\label{sec:rest}

\begin{prop}[reduction]\label{prop:red}
Sign-Schur as a rank-1 / checkerboard / M-matrix fact is
REFUTED as a closing deduction.
The Chebyshev-CD signs on the cosine grid are Dirichlet-zonal
(SATZ); $K^\mu[\Xi]$ inherits those signs in the large
couplings (census, not a SATZ) and not as a bipartition.
No source-defined conjugation bound improves the r382 entry:
the checkerboard positive-remainder majorant crosses $1$ at
$k=\lfloor 2N/5\rfloor$, and the oracle bound stays above $1$
through the wall.
Every $k^*>2N/5$ would have improved
\texttt{adaptive\_band\_from\_entry} mechanically; this round
does not produce one.
The Assist contraction remains the signed $\mu$-OP Gram of
r387/r389, now with a named negative: it is \emph{not} a Schur
bipartition of the $\nu$-nodes.
The hydra-legal remainder is the Dirichlet/Abel (Weyl) energy
of that Gram --- the r389 rest --- not a sign vector.
$F_\varepsilon$ (bound on $\Delta^2\log\tau$ under $F_1$) and
$W_\varepsilon$ are untouched.
No $m_4$.  No RH claim.
\end{prop}

\section{Machine checks}\label{sec:machine}

Numbered lemmas: \texttt{verify\_sign\_schur.py} (same
directory).
Standalone: Dirichlet sign and envelope; Chebyshev-CD;
$2\times 2$ Z-matrix; bookkeeping; checkerboard rank-1;
mesh step.
Construction: FRAME-A Assist, Chebyshev/checkerboard
mass-agreements, $\lvert E\rvert$ mutant, two-period in-phase
and too-late anti-checkerboard, scramble envelope.
Discovery census: \texttt{sign\_schur\_probe.py}
(\texttt{--smoke} / full), core-$42$.
Exit line: \texttt{SIGN SCHUR VERIFIED}.

\section*{Claim boundary}

Finite identities for the Chebyshev--Dirichlet kernel on a
cosine grid, a named $2\times 2$ counterexample to the hoped
M-matrix bound, and named two-period / scramble / unsigned
breaks of a rank-1 sign conjugation.
No statement about zeros of $\zeta(s)$, in either direction.
No RH claim.  No $L^*$ claim.

\begin{thebibliography}{9}
\bibitem{r392}
S.~Hamann,
\textit{The deletion transform, reduced},
standalone note, August 2026,
\texttt{rh/problem/deletion\_transform.tex}.
\bibitem{r387}
S.~Hamann,
\textit{Coherence assist, reduced},
standalone note, August 2026,
\texttt{rh/problem/coherence\_assist.tex}.
\bibitem{r389}
S.~Hamann,
\textit{The Weyl energy of cancellation, reduced},
standalone note, August 2026,
\texttt{rh/problem/weyl\_energy.tex}.
\bibitem{r388}
S.~Hamann,
\textit{The $\Delta$-deformation, reduced},
standalone note, August 2026,
\texttt{rh/problem/delta\_deformation.tex}.
\bibitem{pivot}
S.~Hamann,
\textit{The pivot-band entry lemma, reduced},
standalone note, August 2026,
\texttt{rh/problem/pivot\_entry\_lemma.tex}.
\end{thebibliography}

\end{document}
